\section{Bài tập 2}

Cho mạch điện như hình bên dưới. Hãy sắp xếp (vẽ) lại mạch điện để làm rõ cấu trúc nối tiếp và/hoặc song song của mạch. Sau đó, vận dụng kiến thức đã học để tìm giá trị điện trở tương đương giữa hai đầu $A$ và $F$. Cuối cùng, thực hiện mô phỏng để kiểm tra xem dòng điện qua toàn mạch đã được tính đúng hay chưa.

\given{ex2_given}{0.6\textwidth}{Tìm điện trở tương đương giữa hai đầu $A$ và $F$}

\subsection{Vẽ lại mạch điện}
\drawfig{ex2_redraw}{Mạch điện sau khi vẽ lại dạng hình thang}

\subsection{Tính toán}
\noindent Dây nối giữa nút $E$ và $D$ là dây dẫn lý tưởng nên $E$ và $D$ cùng một nút ($V_E = V_D$).
\begin{itemize}[leftmargin=2em]
  \item Ba điện trở $R_3, R_4, R_5$ nằm trên cùng một nhánh nên mắc nối tiếp:
        \[
        R_{345} = R_3 + R_4 + R_5 = 4\,\Omega + 5\,\Omega + 3\,\Omega = 12\,\Omega
        \]
  \item Giữa hai nút $C$ và $D$, điện trở $R_6$ mắc song song với $R_{345}$:
        \[
        R_{CD\_3456} = R_6 \parallel R_{345} = \frac{R_6 \cdot R_{345}}{R_6 + R_{345}} = \frac{4 \cdot 12}{4 + 12} = 3\,\Omega
        \]
  \item Giữa hai nút $B$ và $E$, nhánh $(R_2 + R_{CD\_3456})$ mắc song song với $R_7$:
        \[
        R_{BE} = R_7 \parallel (R_2 + R_{CD\_3456}) = \frac{6 \cdot (3 + 3)}{6 + (3 + 3)} = \frac{36}{12} = 3\,\Omega
        \]
  \item Giữa hai đầu $A$ và $F$, các điện trở $R_1, R_{BE}$ và $R_8$ mắc nối tiếp:
        \[
        R_{AF} = R_1 + R_{BE} + R_8 = 2\,\Omega + 3\,\Omega + 1\,\Omega = 6\,\Omega
        \]
  \item Dòng điện qua toàn mạch (chạy từ $A$ sang $B$ qua $R_1$):
        \[
        I_{AB} = \frac{V_1}{R_{AF}} = \frac{12\text{ V}}{6\,\Omega} = 2\text{ A}
        \]
\end{itemize}

\begin{dapsobox}
  \begin{itemize}[leftmargin=1.5em]
    \item $R_{CD\_3456} = 3\,\Omega$
    \item $R_{BE} = 3\,\Omega$
    \item $R_{AF} = 6\,\Omega$
    \item $I_{AB} = 2\text{ A}$
  \end{itemize}
\end{dapsobox}

\subsection{Mô phỏng}
\simresult
\simfig{ex2_sim}{Kết quả mô phỏng PSpice của Bài tập 2}
